\documentclass[11pt]{article}
\usepackage[utf8]{inputenc}
\usepackage[margin=1in]{geometry}
\usepackage{amsmath,amssymb}
\usepackage{booktabs}
\usepackage{hyperref}
\usepackage{setspace}
\usepackage{float}
\usepackage{enumitem}
\onehalfspacing

\hypersetup{colorlinks=true,linkcolor=blue!70!black,citecolor=blue!70!black,urlcolor=blue!70!black}

\title{\textbf{VolEdge: Measuring and Trading the\\Variance Risk Premium with Model-Free\\Risk-Neutral Moments}}
\author{Tanishk Yadav\\[2pt]\textit{MS Financial Engineering}\\NYU Tandon School of Engineering\\New York, NY, USA}
\date{Technical note - working draft}

\begin{document}
\maketitle

\begin{abstract}
\noindent
VolEdge is a full-stack research platform built around a single thesis: the gap
between what options \emph{price} and what the underlying \emph{realizes} is the
risk premium, and that gap can be measured and traded. The platform extracts the
risk-neutral distribution model-free via Bakshi-Kapadia-Madan (BKM) moments and
compares it, horizon-matched, against the physical distribution estimated from
realized returns; the spread $Q-P$ is surfaced for variance, skew, and kurtosis.
A no-look-ahead walk-forward engine then backtests premium-driven strategies.
We report two honest results on free data (2019-2024, net of modeled costs):
a GARCH(1,1) volatility-targeted strategy beat SPY buy-and-hold on a risk-adjusted
basis (Sharpe 0.85 vs.\ 0.68) with 41\% lower maximum drawdown, while a naive
long-only variance-premium harvest \emph{lost} money (Sharpe $-0.16$), empirically
confirming the premium as compensation for crash risk rather than harvestable
alpha. The BKM extractor is validated against Black-Scholes chains of known
volatility, and the backtest engine reproduces buy-and-hold to $\pm 0.000$ pp at
any rebalance frequency.

\medskip
\noindent\textbf{Keywords:} variance risk premium, risk-neutral moments, BKM,
volatility targeting, walk-forward backtesting, options.
\end{abstract}

\section{Introduction}

In a frictionless world the option-implied (risk-neutral, $Q$) distribution would
equal the physical (real-world, $P$) distribution. It does not: $Q$ embeds a risk
premium, so an option's price is \emph{not} a probability. The wedge between the
two measures is the economic object of interest. VolEdge is organized entirely
around keeping $Q$ and $P$ as distinct, separately estimated quantities and
studying their difference.

The platform measures the premium live for any ticker, and - because a historical
per-name option surface is not available on free data - backtests the idea at the
index level using the VIX as the risk-neutral volatility proxy. The two
measurements are deliberately distinct; the backtest's role is to test whether the
premium the live tool surfaces is economically real and harvestable, not to
reproduce BKM historically.

\section{Methods}

\subsection{Risk-neutral moments (BKM)}
Following Bakshi, Kapadia, and Madan (2003), any twice-differentiable payoff is
spanned by a continuum of out-of-the-money (OTM) options. The variance, cubic, and
quartic contracts therefore have known strike weights, and their prices
$V,W,X$ are obtained by integrating OTM call and put prices across strikes. With
$\log$-moneyness $\ell=\ln(K/S)$ the call and put integrands unify:
\[
V=\!\int\! \frac{2(1-\ell)}{K^2}O(K)\,dK,\quad
W=\!\int\! \frac{6\ell-3\ell^2}{K^2}O(K)\,dK,\quad
X=\!\int\! \frac{12\ell^2-4\ell^3}{K^2}O(K)\,dK,
\]
where $O(K)$ is the OTM option price. Risk-neutral variance, skewness, and
(excess) kurtosis follow in closed form, with a single $e^{rT}$ entering the
moment formulas (not the integrands). Integration uses composite Simpson's rule
over the irregular strike grid with a trapezoidal fallback. No distributional
assumption is made.

\subsection{Physical moments and the premium}
The physical distribution is estimated horizon-matched: because skew and kurtosis
are horizon-dependent, they are computed on \emph{overlapping $T$-day log returns}
rather than daily returns, where $T$ matches the BKM tenor. The premium is the
difference at each tenor,
\[
\text{VRP}=\sigma^2_Q-\sigma^2_P,\qquad
\text{skew premium}=s_Q-s_P,\qquad
\text{kurt premium}=k_Q-k_P .
\]
A persistently positive variance premium means options price more variance than
realizes; a more negative $Q$-skew than $P$-skew means the market pays up for
crash protection beyond what history delivered.

\subsection{Mean-reversion diagnostics}
As a descriptive layer, the platform reports three independent reversion estimators
on price, log-price, and realized volatility: the Ornstein-Uhlenbeck half-life
from the AR(1) speed ($-\ln 2/\ln\phi$), the Hurst exponent, and a
Dickey-Fuller-style $t$-statistic, plus a rolling half-life. A series is labelled
mean-reverting only when a finite in-sample half-life corroborates a Hurst below
$0.5$.

\subsection{Walk-forward engine}
Strategies are evaluated day-by-day; at each decision point the strategy sees only
data up to the prior close. Strategy code is AST-validated and sandboxed
(whitelisted numerical imports, no I/O or network). Positions are marked to
market \emph{before} any rebalance, so returns accrue correctly at every rebalance
frequency; turnover cost is charged on the actual drift-to-target trade with a
per-leg model (spread, impact, slippage, commission).

\section{Validation}

\paragraph{BKM.} On a Black-Scholes chain of known 20\% volatility the extractor
recovers $\sigma_Q\approx0.20$ with skewness and excess kurtosis near zero, and
converges toward truth as strike coverage widens; the residual bias is strike
truncation, an inherent limitation of BKM on any finite chain. A synthetic
put-skew chain correctly yields negative risk-neutral skew. On live SPY chains
(via free Yahoo Finance data) the moments show a sensible term structure
(e.g.\ 30d skew $\approx -2.2$, kurtosis $\approx 9$; 60d $\approx -1.6$, $3.4$),
consistent with published equity-index estimates.

\paragraph{Engine.} A 100\%-SPY strategy reproduces SPY buy-and-hold to
$\pm 0.000$ percentage points at rebalance frequencies of 1, 5, 21, and 63 days,
confirming the mark-to-market is correct and frequency-independent.

\section{Results}

All figures are out-of-sample, on free Yahoo Finance data, net of modeled
transaction costs.

\begin{table}[H]
\centering
\begin{tabular}{lrrrrr}
\toprule
Strategy & Return & Ann. & Vol & Sharpe & Max DD\\
\midrule
GARCH vol-managed SPY (2020-24) & $+75.1\%$ & $11.9\%$ & $13.9\%$ & $\mathbf{0.85}$ & $\mathbf{-19.7\%}$\\
\quad SPY buy-and-hold (same)     & $+94.6\%$ & $14.3\%$ & $21.0\%$ & $0.68$ & $-33.7\%$\\
VIX variance-premium harvest (2019-24) & $-15.4\%$ & $-2.9\%$ & $17.8\%$ & $-0.16$ & $-44.0\%$\\
\quad SPY (2019-24)              & $+123.9\%$ & $15.1\%$ & $20.1\%$ & $0.75$ & $-33.7\%$\\
\bottomrule
\end{tabular}
\caption{Backtested premium-driven strategies vs.\ SPY.}
\end{table}

\paragraph{GARCH volatility targeting (the positive result).}
A GARCH(1,1) conditional-variance forecast scales SPY exposure to a constant
15\% volatility target (Moreira-Muir), the remainder in T-bills. On an
apples-to-apples basis it delivered a higher Sharpe (0.85 vs.\ 0.68, $+25\%$) at
41\% lower maximum drawdown and a realized 13.9\% volatility. Because volatility is
forecastable while expected returns are not, scaling down in high-forecast-vol
states improves the realized risk/return trade-off - and the variance premium
makes those states expensive to be long, so the de-risking is compensated.

\paragraph{Variance-premium harvest (the reality check).}
A long-only short-volatility harvester (short via SVXY when $\text{VIX}^2-$realized
is positive, scaling down as VIX rises and stepping out above VIX 30) \emph{lost}
money over 2019-2024 net of costs. This is the intended finding, not a hidden
failure: across the COVID-2020 and 2022 volatility spikes the premium behaves as
compensation for crash risk, not harvestable alpha with simple long-only
instruments - which is precisely why the premium exists. The platform thus both
measures the premium and demonstrates that harvesting it naively is dangerous.

\section{Limitations}
Yahoo Finance quotes are $\sim$15-minute-delayed snapshots with no SLA; deep-OTM
strikes can be stale (mitigated by a liquidity filter and a $\pm 20\%$ strike
window). Kurtosis is the least reliable BKM moment on free data. The live premium
(per-name BKM snapshot) and the backtest (index VIX proxy) measure the idea two
different ways. The walk-forward accounting is weight-based rather than a full
position ledger. None of the reported numbers are tuned via parameter search behind
the headline.

\section{Conclusion}
VolEdge closes the loop from measurement to execution for the variance risk
premium: model-free risk-neutral moments versus horizon-matched realized moments,
validated extractors, and a frequency-correct walk-forward engine. The empirical
takeaway is that the premium is easy to measure and hard to keep - disciplined
volatility targeting is compensated, while naive premium harvesting is not.

\paragraph{Code.} \url{https://github.com/tanishhky/voledge}

\end{document}
